\chapter{Gastrectomy}

\section*{Introduction \& Clinical Indications}
Gastrectomy is a major surgical procedure involving the partial or complete removal of the stomach, a vital organ in the upper gastrointestinal tract responsible for initial food digestion. This intervention is primarily undertaken to treat malignant conditions such as gastric adenocarcinoma, but it is also indicated for certain severe benign diseases that are refractory to medical management. These benign indications include intractable peptic ulcer disease with complications, severe gastric polyps with malignant potential, or gastrointestinal stromal tumors (GIST). The extent of the resection, whether subtotal (partial) or total, depends critically on the underlying pathology, its precise location within the stomach, and its stage. Following the removal of stomach tissue, meticulous surgical reconstruction of the gastrointestinal tract is essential to restore alimentary continuity, most commonly utilizing techniques like Billroth I, Billroth II, or Roux-en-Y anastomoses. This complex procedure carries significant clinical implications for patient prognosis, particularly in the context of cancer treatment, and profoundly impacts long-term quality of life due to altered digestive physiology.

\begin{itemize}
  \item Stomach resection
  \item Gastric resection
  \item Subtotal gastrectomy
  \item Total gastrectomy
  \item Partial gastrectomy
  \item Gastric removal surgery
\end{itemize}

The fundamental principle of gastrectomy is the surgical eradication of diseased gastric tissue while preserving as much healthy tissue and function as possible. In cases of gastric malignancy, the primary goal is to achieve a complete oncologic resection (R0 resection), meaning the removal of all visible and microscopic tumor with clear surgical margins. This often necessitates an extensive lymphadenectomy, removing regional lymph nodes that may harbor metastatic cancer cells, thereby improving staging accuracy and reducing recurrence risk. The rationale for this aggressive approach in cancer is to maximize the chance of cure and improve long-term survival.

For benign conditions, the principle focuses on removing the source of pathology, such as a bleeding ulcer, a symptomatic tumor, or a focus of high-grade dysplasia, to alleviate symptoms and prevent life-threatening complications like perforation, hemorrhage, or malignant transformation. Following resection, the digestive tract must be meticulously reconstructed to allow for the passage of food, bile, and pancreatic secretions, aiming to minimize post-gastrectomy syndromes and maintain nutritional status. The choice of reconstruction method (e.g., Billroth I, Billroth II, Roux-en-Y) is guided by the extent of resection, the need to prevent reflux, and the surgeon's preference, each designed to re-establish continuity between the esophagus or remaining stomach and the small intestine in a physiologically sound manner.

The history of gastrectomy is marked by significant advancements in surgical technique and understanding of gastric physiology, transforming a once-lethal procedure into a life-saving intervention. The pioneering work began in the late 19th century. Theodor Billroth, an Austrian surgeon, performed the first successful partial gastrectomy for gastric cancer in 1881, a landmark operation that demonstrated the feasibility of gastric resection and introduced the Billroth I reconstruction, where the remaining stomach is anastomosed directly to the duodenum. This initial success paved the way for further innovation.

Shortly thereafter, in 1885, Billroth introduced the Billroth II reconstruction, which involves closing the duodenal stump and anastomosing the remaining stomach to the jejunum, a technique often favored when a tension-free Billroth I anastomosis is not possible or when duodenal pathology precludes its use. Further refinement came with C\'{e}sar Roux's introduction of the Roux-en-Y gastrojejunostomy in 1893, initially for benign conditions, but later widely adopted for total gastrectomy and complex reconstructions due to its ability to divert bile and pancreatic secretions away from the gastric remnant, thereby reducing reflux and associated complications. The 20th century saw continuous improvements in anesthesia, critical care, surgical instruments, and understanding of fluid and electrolyte balance, making gastrectomy safer and more effective. More recently, the advent of minimally invasive laparoscopic and robotic techniques has further refined the procedure, offering potential benefits such as reduced pain, shorter hospital stays, and faster recovery for selected patients, particularly for distal gastrectomies.

Gastrectomy is indicated for a variety of conditions, predominantly malignant, but also for severe benign pathologies.

\begin{itemize}
  \item To treat gastric adenocarcinoma, aiming for curative resection with clear margins and regional lymphadenectomy, often guided by preoperative staging (e.g., TNM staging).
  \item To remove gastrointestinal stromal tumors (GIST) of the stomach, especially those larger than 2 cm, rapidly growing, or with high-risk features identified on biopsy.
  \item To manage severe, refractory peptic ulcer disease complicated by perforation, intractable bleeding, or gastric outlet obstruction that has failed extensive medical therapy and endoscopic interventions.
  \item To resect high-grade dysplasia or early gastric cancer (e.g., T1a or T1b) identified during endoscopic surveillance, preventing progression to invasive malignancy.
  \item For palliative management of advanced gastric cancer causing intractable bleeding, obstruction, or severe pain, even in the presence of metastatic disease, to improve quality of life.
  \item To remove large, symptomatic benign gastric polyps (e.g., adenomatous polyps) with significant malignant potential or those causing bleeding or obstruction.
  \item In rare cases, to address severe, irreparable gastric trauma (e.g., extensive gunshot wounds) or caustic injury that has resulted in necrosis or perforation.
  \item For complications of Zollinger-Ellison syndrome, such as severe, recurrent ulceration or perforation, when medical management with proton pump inhibitors is insufficient or poorly tolerated.
\end{itemize}

Gastrectomy serves as a primary treatment in most cases, particularly for localized gastric adenocarcinoma where it offers the best chance for cure. For early-stage gastric cancer without evidence of distant metastasis, a gastrectomy with adequate lymphadenectomy is considered the cornerstone of curative therapy, often combined with neoadjuvant (preoperative) or adjuvant (postoperative) chemotherapy and/or radiation to improve outcomes. Similarly, for severe benign conditions like intractable bleeding ulcers, large symptomatic GISTs, or high-grade dysplasia, gastrectomy is a primary intervention when medical or endoscopic treatments have failed or are deemed insufficient to control the disease or prevent complications.

However, gastrectomy can also function as a secondary or palliative procedure. In advanced gastric cancer with widespread metastasis, a gastrectomy may be performed secondarily to alleviate severe symptoms such as intractable bleeding, gastric outlet obstruction, or severe pain, thereby improving the patient's quality of life even if a cure is not achievable. This is often referred to as a palliative gastrectomy. It may also be a salvage procedure for complications arising from previous gastric surgeries, such as recurrent anastomotic strictures or fistulas, or for treatment failures of other modalities. The decision for gastrectomy, therefore, is always carefully weighed against the patient's overall health, disease stage, and potential for benefit.

\section*{Relevant Surgical Anatomy}
The stomach is a J-shaped organ situated in the upper abdomen, primarily in the epigastric and left hypochondriac regions. It is divided into the cardia, fundus, body, antrum, and pylorus. Its walls consist of four layers: mucosa, submucosa, muscularis propria (with inner oblique, middle circular, and outer longitudinal layers), and serosa. Key anatomical landmarks include the greater curvature (lateral, convex border) and lesser curvature (medial, concave border), which serve as attachments for the greater and lesser omenta, respectively. The lesser omentum connects the lesser curvature to the liver, containing the hepatogastric and hepatoduodenal ligaments, with the latter forming the anterior boundary of the epiploic foramen (of Winslow) and housing the portal triad (hepatic artery, portal vein, common bile duct).

The arterial supply to the stomach is rich, originating primarily from the celiac axis. The lesser curvature is supplied by the left and right gastric arteries (branches of the celiac and common hepatic arteries, respectively). The greater curvature receives blood from the left and right gastroepiploic arteries (branches of the splenic and gastroduodenal arteries, respectively). The fundus is supplied by the short gastric arteries, originating from the splenic artery. Venous drainage largely parallels the arterial supply, emptying into the portal venous system. Lymphatic drainage is extensive and follows the arterial supply, with primary nodal basins along the greater and lesser curvatures (perigastric nodes), the celiac axis, and the splenic hilum, crucial for oncologic staging. The vagus nerves (anterior and posterior trunks) provide parasympathetic innervation, controlling gastric motility and acid secretion, while sympathetic innervation arises from the celiac plexus. Surgical landmarks include the gastroesophageal junction proximally and the pylorus distally, which marks the transition to the duodenum. The relationship of the stomach to adjacent organs like the spleen, pancreas, liver, and diaphragm is critical during mobilization and lymphadenectomy to avoid iatrogenic injury.

\section*{Preoperative Workup \& Preparation}
Thorough preoperative preparation is crucial to optimize patient outcomes and minimize complications associated with gastrectomy. Patients are typically required to be NPO (nil per os) for at least 6-8 hours prior to surgery to reduce the risk of aspiration during anesthesia.

Key preparatory steps include:
\begin{itemize}
  \item \textbf{Medical Evaluation:} A comprehensive medical history and physical examination are performed. This includes a detailed assessment of cardiac, pulmonary, and renal function, often requiring consultations with cardiology or pulmonology specialists to optimize chronic conditions.
  \item \textbf{Laboratory Tests:} Baseline blood tests are ordered, including a complete blood count (CBC) to check for anemia, a comprehensive metabolic panel (electrolytes, liver function tests, renal function tests), coagulation profile (PT/INR, PTT) to assess bleeding risk, and blood type and crossmatch for potential transfusions.
  \item \textbf{Imaging and Endoscopy:} Preoperative imaging such as CT scans of the chest, abdomen, and pelvis is essential for accurate staging of gastric cancer and assessing resectability. Upper endoscopy with biopsy is typically performed to confirm diagnosis, delineate tumor extent, and rule out multifocal disease. PET scans may be used for further staging.
  \item \textbf{Nutritional Optimization:} Many patients undergoing gastrectomy, especially for cancer, are malnourished. Nutritional assessment (e.g., albumin, prealbumin levels, weight loss history) and optimization, potentially including preoperative nutritional support (oral supplements, enteral feeding via jejunostomy tube, or total parenteral nutrition), are vital to improve surgical tolerance and wound healing. Anemia, if present, should be corrected.
  \item \textbf{Medication Review:} All current medications are reviewed. Anticoagulants (e.g., warfarin, DOACs) may need to be stopped or bridged with heparin, and diabetic medications adjusted to prevent hypo- or hyperglycemia.
  \item \textbf{Bowel Preparation:} While not always routinely required for gastrectomy, some surgeons may recommend a mechanical bowel preparation, especially if a colonic resection is anticipated or if the reconstruction involves a long segment of jejunum.
  \item \textbf{Antibiotic Prophylaxis:} Intravenous broad-spectrum antibiotics (e.g., cefazolin plus metronidazole) are administered within one hour prior to incision to reduce the risk of surgical site infection.
  \item \textbf{Informed Consent:} A detailed discussion of the procedure, potential risks, benefits, alternatives, and expected postoperative course is conducted, ensuring the patient provides fully informed consent.
\end{itemize}

Gastrectomy is a major surgery, and careful consideration of patient factors is paramount.

\textbf{Situations requiring staging, optimization, or an alternative treatment plan:}
\begin{itemize}
  \item Inability to tolerate general anesthesia due to severe, uncorrectable medical comorbidities (e.g., severe cardiac failure, uncontrolled respiratory failure).
  \item Widespread, unresectable metastatic disease where the primary tumor is asymptomatic and gastrectomy would not offer any palliative benefit, only increased morbidity.
  \item Uncontrolled severe coagulopathy that cannot be corrected preoperatively, posing an unacceptably high risk of hemorrhage.
  \item Acute, severe myocardial infarction or unstable angina, requiring cardiac stabilization before any elective surgery.
\end{itemize}

\textbf{Relative Contraindications \& Precautions:}
\begin{itemize}
  \item Severe cardiac or pulmonary disease that significantly increases operative risk, requiring careful optimization, multidisciplinary team assessment, and a thorough risk-benefit discussion with the patient.
  \item Extreme frailty or advanced age, where the physiological stress of surgery may outweigh potential benefits, necessitating a comprehensive geriatric assessment.
  \item Severe malnutrition or cachexia, which significantly impairs wound healing and increases complication rates; aggressive preoperative nutritional support is essential to improve surgical tolerance.
  \item Portal hypertension with significant esophageal or gastric varices, which substantially increases the risk of intraoperative and postoperative bleeding, requiring careful management.
  \item Active infection at a distant site (e.g., pneumonia, urinary tract infection), which should be treated and resolved before elective surgery to prevent systemic complications.
  \item Extensive local invasion of the tumor into vital unresectable structures (e.g., aorta, superior mesenteric artery), making an R0 resection impossible without significant, unacceptable morbidity or mortality.
  \item Prior extensive abdominal surgery leading to dense adhesions, which can increase operative time and risk of iatrogenic injury.
\end{itemize}

\section*{Surgical Technique \& Procedure}
Gastrectomy is performed under general anesthesia, typically with the patient in a supine position, often with a slight reverse Trendelenburg tilt to improve exposure of the upper abdomen. The procedure can be carried out via an open laparotomy incision (e.g., upper midline or bilateral subcostal) or through a minimally invasive approach using laparoscopic or robotic techniques, which may involve several small incisions.

The key operative steps generally include:
\begin{enumerate}
  \item \textbf{Exploration and Assessment:} After gaining abdominal access, a thorough exploration of the abdominal cavity is performed to confirm the extent of the disease, rule out unforeseen metastases (e.g., peritoneal carcinomatosis, liver metastases), and assess the resectability of the primary tumor. Peritoneal washings for cytology may be taken in cancer cases.
  \item \textbf{Mobilization of the Stomach:} The stomach is carefully mobilized from its attachments. This involves dividing the gastrocolic ligament to enter the lesser sac, ligating the short gastric vessels along the greater curvature, and dissecting the lesser omentum along the lesser curvature, while meticulously preserving the spleen if possible.
  \item \textbf{Lymphadenectomy:} For gastric cancer, a meticulous dissection of regional lymph nodes is performed according to oncologic guidelines (e.g., D1 or D2 lymphadenectomy). This involves removing lymph nodes along the greater and lesser curvatures, around the celiac axis, hepatic artery, splenic artery, and other perigastric stations, depending on the tumor location and extent.
  \item \textbf{Resection:} The diseased portion of the stomach is resected. For a subtotal gastrectomy, the distal stomach is removed, typically including the antrum and a portion of the body. For a total gastrectomy, the entire stomach is excised, including the gastroesophageal junction proximally and the pylorus distally. Surgical margins are carefully assessed intraoperatively (e.g., frozen section analysis) to ensure clear resection.
  \item \textbf{Reconstruction of Alimentary Continuity:} After resection, the digestive tract is restored. Common methods include:
    \begin{itemize}
      \item \textbf{Billroth I:} The remaining stomach (after subtotal gastrectomy) is anastomosed directly to the duodenum.
      \item \textbf{Billroth II:} The duodenal stump is closed, and the remaining stomach is anastomosed to a loop of jejunum (gastrojejunostomy).
      \item \textbf{Roux-en-Y:} A loop of jejunum is brought up and anastomosed to the remaining stomach or esophagus (after total gastrectomy), with a separate enteroenterostomy further down to divert bile and pancreatic secretions away from the gastroesophageal or gastrojejunal anastomosis.
    \end{itemize}
  \item \textbf{Hemostasis and Drain Placement:} Meticulous hemostasis is achieved throughout the surgical field. A surgical drain (e.g., Jackson-Pratt) may be placed near the anastomosis or in the lesser sac to monitor for leaks or fluid collections.
  \item \textbf{Closure:} The abdominal incision is closed in layers, or laparoscopic port sites are closed, ensuring proper fascial and skin approximation.
\end{enumerate}

\section*{Complications \& Risks}
Gastrectomy is a major surgical procedure associated with a range of potential risks and complications, which can be broadly categorized as general surgical risks and procedure-specific risks.

\begin{itemize}
  \item \textbf{General Surgical Risks:}
    \begin{itemize}
      \item \textbf{Bleeding:} Intraoperative or postoperative hemorrhage, potentially requiring blood transfusions or re-operation to control the source.
      \item \textbf{Infection:} Surgical site infection (SSI) at the incision, intra-abdominal abscess formation, or systemic infections like pneumonia or urinary tract infection.
      \item \textbf{Anesthesia Complications:} Adverse reactions to anesthetic agents, respiratory depression, cardiac events (e.g., myocardial infarction, arrhythmias), or aspiration pneumonia.
      \item \textbf{Deep Vein Thrombosis (DVT) and Pulmonary Embolism (PE):} Formation of blood clots in the deep veins of the legs that can dislodge and travel to the lungs, a potentially fatal complication.
      \item \textbf{Cardiac Events:} Myocardial infarction, arrhythmias, or exacerbation of heart failure, particularly in patients with pre-existing cardiac disease.
      \item \textbf{Stroke:} Cerebrovascular accident due to perioperative hemodynamic changes or emboli.
      \item \textbf{Wound Complications:} Seroma, hematoma, wound dehiscence, or incisional hernia.
    \end{itemize}
  \item \textbf{Procedure-Specific Risks:}
    \begin{itemize}
      \item \textbf{Anastomotic Leak:} Leakage of digestive contents from the surgical connection between the remaining stomach/esophagus and intestine, a serious and life-threatening complication requiring urgent re-intervention.
      \item \textbf{Dumping Syndrome:} Rapid emptying of undigested food into the small intestine, leading to early symptoms (nausea, vomiting, diarrhea, abdominal cramps, lightheadedness) or late symptoms (hypoglycemia) particularly after meals high in sugar.
      \item \textbf{Nutritional Deficiencies:} Long-term malabsorption leading to deficiencies in vitamin B12 (requiring lifelong intramuscular supplementation after total gastrectomy due to loss of intrinsic factor), iron, calcium, vitamin D, and fat-soluble vitamins.
      \item \textbf{Marginal Ulceration:} Ulcers forming at the anastomotic site, especially common after Billroth II reconstruction, often requiring medical management or re-operation.
      \item \textbf{Gastric Stasis or Delayed Gastric Emptying:} Impaired motility of the remaining stomach or small bowel, leading to prolonged nausea, vomiting, and inability to tolerate oral intake.
      \item \textbf{Reflux Esophagitis/Gastritis:} Bile or intestinal fluid reflux into the esophagus or gastric remnant, causing inflammation, pain, and discomfort, particularly after Billroth I or II reconstructions.
      \item \textbf{Afferent Loop Syndrome:} Obstruction of the afferent limb (the segment of jejunum leading to the anastomosis in Billroth II), causing pain, nausea, vomiting, and potentially cholangitis or pancreatitis.
      \item \textbf{Internal Hernia:} Protrusion of bowel through a defect created during reconstruction, particularly with Roux-en-Y anastomoses, leading to bowel obstruction.
      \item \textbf{Anastomotic Stricture:} Narrowing at the site of the surgical connection, potentially requiring endoscopic dilation or surgical revision.
      \item \textbf{Pancreatitis:} Inflammation of the pancreas due to surgical manipulation, injury to the pancreatic duct, or obstruction of the pancreaticobiliary outflow.
    \end{itemize}
\end{itemize}

\section*{Postoperative Care \& Recovery}
Immediately after gastrectomy, the patient is transferred to the Post-Anesthesia Care Unit (PACU) for close monitoring of vital signs, pain levels, and potential signs of immediate complications such as bleeding or respiratory distress. Pain management is initiated, often with patient-controlled analgesia (PCA) or epidural catheters. Intravenous fluids are administered to maintain hydration, and a nasogastric tube may be in place to decompress the stomach and prevent nausea, though its routine use is debated and often removed early.

Over the first 24-48 hours, the patient is encouraged to ambulate early to prevent complications like DVT and pneumonia. Oral intake is typically restricted (NPO) initially, with gradual progression from clear liquids to a soft, low-fat, high-protein, small-volume diet as bowel function returns and tolerance improves. Close monitoring for signs of anastomotic leak (e.g., persistent tachycardia, fever, unexplained abdominal pain, purulent drain output, elevated inflammatory markers) is paramount. Surgical drains are usually removed when output is minimal and serous. Patients are educated on dietary modifications to manage dumping syndrome and other post-gastrectomy symptoms. The typical hospital stay ranges from 5 to 10 days, depending on the extent of surgery, the patient's recovery trajectory, and the absence of complications. Long-term recovery involves adherence to specific dietary guidelines, often involving small, frequent meals, and may require lifelong vitamin B12 supplementation, especially after total gastrectomy, due to the loss of intrinsic factor production. Full recovery and return to normal activities can take several weeks to months.

During gastrectomy, the surgeon meticulously documents intraoperative findings, including the precise location and size of the tumor or pathology, its relationship to adjacent organs, the presence of any suspicious lymph nodes, and the extent of any metastatic disease. This information is crucial for accurate staging and guiding subsequent treatment. The resected stomach specimen, along with all removed lymph nodes, is then sent to pathology for detailed microscopic examination.

The pathology report on the resected tissues represents the definitive diagnosis and provides critical information. For cancer patients, it will detail the specific type of cancer (e.g., adenocarcinoma, GIST), its histological grade, the depth of tumor invasion into the stomach wall (T-stage), and most importantly, the status of the surgical margins (R0 indicates clear margins, R1 indicates microscopic positive margins, R2 indicates macroscopic positive margins). The number of lymph nodes examined and the number found to contain cancer cells are crucial for nodal staging (N-stage). Other important features include lymphovascular and perineural invasion, which are prognostic indicators. For benign conditions, the pathology report confirms the diagnosis (e.g., peptic ulcer, benign polyp, leiomyoma) and ensures no unexpected malignancy was present. A clear (R0) margin in cancer cases is the most favorable result, indicating complete tumor removal.

A normal, successful postoperative course following gastrectomy involves a predictable progression towards recovery and functional independence. Patients can expect initial pain, which should be well-managed with analgesia and gradually decrease over the first few days. Bowel function typically returns within 3-5 days, marked by the passage of flatus and then bowel movements, allowing for a gradual advancement of diet. The surgical wound should heal progressively, with sutures or staples typically removed around 10-14 days post-op. Patients should experience a resolution of their preoperative symptoms (e.g., bleeding, obstruction, intractable pain from an ulcer or tumor). While some degree of dietary adjustment is almost always necessary, a successful course means the patient learns to manage these changes effectively, minimizing symptoms like dumping syndrome or early satiety. Energy levels will gradually improve, and the patient should be able to resume light activities within 2-4 weeks, with a progressive return to normal daily routines and work over 2-3 months. The ultimate goal is a return to a good quality of life with stable weight and nutritional status.

Post-gastrectomy follow-up is essential for monitoring recovery, managing long-term complications, and, for cancer patients, surveillance for recurrence.

\begin{itemize}
  \item \textbf{Surgical Follow-up:} Initial postoperative visits with the surgeon typically occur 2-4 weeks after discharge to assess wound healing, discuss pathology results, and address any immediate concerns.
  \item \textbf{Nutritional Counseling and Supplementation:} Lifelong nutritional counseling is crucial. Patients who undergo total gastrectomy will require lifelong intramuscular vitamin B12 injections due to the loss of intrinsic factor. Iron, calcium, and vitamin D levels are monitored regularly, and supplements are often prescribed to prevent deficiencies.
  \item \textbf{Oncology Follow-up:} For cancer patients, regular follow-up with an oncologist is mandatory to plan and administer adjuvant therapies (chemotherapy, radiation) and for long-term surveillance for recurrence, which may include periodic CT scans, endoscopic evaluations, and tumor marker monitoring.
  \item \textbf{Laboratory Monitoring:} Regular blood tests, including CBC, iron studies, vitamin B12 levels, and calcium/vitamin D, are performed every 3-6 months initially, then annually, to detect and manage nutritional deficiencies.
  \item \textbf{Dietary Adjustments:} Patients are advised on specific dietary modifications, such as eating small, frequent meals (6-8 per day), avoiding high-sugar foods, separating liquids from solids, and chewing thoroughly, to mitigate symptoms of dumping syndrome and promote optimal digestion.
  \item \textbf{Activity Resumption:} Gradual increase in physical activity is encouraged, with avoidance of heavy lifting or strenuous abdominal exercises for several weeks to months, depending on the type of incision and individual healing.
  \item \textbf{Warning Signs:} Patients are educated on warning signs that necessitate immediate medical attention, including persistent fever, severe or worsening abdominal pain, persistent nausea or vomiting, jaundice, dark or bloody stools, significant unexplained weight loss, or difficulty swallowing.
\end{itemize}
Adherence to these follow-up recommendations is critical for optimizing long-term health and detecting potential issues early.

\section*{Outcomes \& Prognosis}
Gastrectomy is most commonly performed for gastric cancer, which is the fifth most common cancer worldwide and the fourth leading cause of cancer-related death. Its incidence shows significant geographical variation, with the highest rates observed in East Asia (e.g., Korea, Japan, China), Eastern Europe, and parts of South America. In Western countries, the incidence has been declining over the past few decades, though it remains a significant health concern, particularly for proximal gastric and gastroesophageal junction cancers. Gastric cancer typically affects older individuals, with a median age at diagnosis in the late 60s, and is more prevalent in men than women. While cancer is the predominant indication, gastrectomy is also performed for benign conditions, albeit less frequently, with peptic ulcer disease being a historical indication that has largely been supplanted by medical therapy. The overall mortality rate for gastrectomy has significantly decreased over the decades due to advancements in surgical techniques, anesthesia, and postoperative care, but it still carries a morbidity rate ranging from 20\% to 40\%, with major complications occurring in 10\% to 20\% of cases, varying by patient factors, surgical complexity, and the extent of lymphadenectomy.

The outcomes and prognosis following gastrectomy are highly dependent on the underlying indication, particularly for gastric cancer. For early-stage gastric cancer (Stage I), the 5-year survival rate after curative gastrectomy with adequate lymphadenectomy can exceed 90\%. However, for more advanced stages, the prognosis significantly worsens, with 5-year survival rates dropping to 20-30\% for Stage III disease, despite surgery and adjuvant therapies. Achieving an R0 resection (clear surgical margins) is the most critical prognostic factor for cancer patients, significantly impacting long-term survival and recurrence rates. Other important prognostic indicators include the number of positive lymph nodes, tumor depth, tumor grade, and the presence of lymphovascular or perineural invasion. For benign conditions, gastrectomy typically offers excellent outcomes in terms of symptom resolution and prevention of complications, with a very low recurrence rate of the original pathology. Functional outcomes, however, can be variable, with many patients experiencing some degree of post-gastrectomy syndrome (e.g., dumping syndrome, nutritional deficiencies) that requires long-term dietary management and supplementation, impacting their quality of life. Landmark clinical trials and meta-analyses, such as those informing the National Comprehensive Cancer Network (NCCN) guidelines, consistently demonstrate the survival benefit of complete surgical resection with adequate lymphadenectomy for localized gastric cancer.

\section*{References}
\begin{enumerate}
  \item MedlinePlus. Gastrectomy. U.S. National Library of Medicine; 2024. Available from: \url{https://medlineplus.gov/gastrectomy.html}
  \item Sabiston Textbook of Surgery: The Biological Basis of Modern Surgical Practice. 22nd ed. Elsevier; 2022.
  \item National Comprehensive Cancer Network (NCCN). NCCN Clinical Practice Guidelines in Oncology: Gastric Cancer. Version 2.2024. Available from: \url{https://www.nccn.org/professionals/physician_gls/pdf/gastric.pdf}
  \item Schwartz's Principles of Surgery. 11th ed. McGraw-Hill Education; 2019.
\end{enumerate}

\clearpage
